\documentclass[11pt]{article}
% ===========================================================================
%  Shared preamble for all MPM2D worksheets — input right after \documentclass.
%  (Files beginning with "_" are skipped by mpm2d-worksheet-publish.mjs.)
% ===========================================================================
\usepackage[letterpaper,top=2.2cm,bottom=1.9cm,left=1.6cm,right=1.6cm,headheight=28pt,footskip=24pt]{geometry}
\usepackage{amsmath,amssymb}
\usepackage[T1]{fontenc}
\usepackage[utf8]{inputenc}
\usepackage{lmodern}
\usepackage{xcolor}
\usepackage[most]{tcolorbox}
\usepackage{enumitem}
\usepackage{multicol}
\usepackage{fancyhdr}
\usepackage{tikz}
\usetikzlibrary{calc}
\usepackage{pgfplots}
\pgfplotsset{compat=1.18}

\definecolor{exblue}{HTML}{4A90E2}
\definecolor{exbg}{HTML}{E6F3FF}
\definecolor{qorange}{HTML}{E69138}
\definecolor{qbg}{HTML}{FFF7CC}
\definecolor{keygray}{HTML}{475569}
\definecolor{keybg}{HTML}{F1F5F9}
\definecolor{iagreen}{HTML}{14653B}
\definecolor{titleblue}{HTML}{1D4ED8}
% Rotating palette so each worked example gets its own colour.
\definecolor{ex0}{HTML}{2563A0}\definecolor{ex1}{HTML}{3B7D3B}\definecolor{ex2}{HTML}{8A5A00}
\definecolor{ex3}{HTML}{A3327A}\definecolor{ex4}{HTML}{6D28A3}\definecolor{ex5}{HTML}{B08900}
\definecolor{ex6}{HTML}{0E7490}\definecolor{ex7}{HTML}{9A3412}\definecolor{ex8}{HTML}{15803D}

\pagestyle{fancy}
\fancyhf{}
\fancyhead[L]{\footnotesize NAME: \underline{\hspace{3.4cm}}}
\fancyhead[C]{\textbf{Integration Academy}\\[-2pt]{\scriptsize Math Simplified \textperiodcentered{} Success Amplified}}
\fancyhead[R]{\footnotesize MPM2D}
\fancyfoot[L]{\scriptsize dr.merie@IntegrationAcademy.ca}
\fancyfoot[C]{\scriptsize\thepage}
\fancyfoot[R]{\scriptsize IntegrationAcademy.ca}
\renewcommand{\headrulewidth}{1.2pt}
\renewcommand{\footrulewidth}{0.4pt}
\setlength{\parindent}{0pt}
\setlength{\parskip}{4pt}

% Examples are NOT breakable, so each example + its solution stays on one page.
\newtcolorbox{example}[2]{colframe=#1,colback=#1!7,coltitle=white,colbacktitle=#1,fonttitle=\bfseries,title={#2},arc=3pt,boxrule=1pt}
\newtcolorbox{qbox}[1]{colframe=qorange,colback=qbg,coltitle=white,colbacktitle=qorange,fonttitle=\bfseries,title={#1},arc=3pt,breakable,boxrule=1pt}
\newtcolorbox{keybox}{colframe=keygray,colback=keybg,coltitle=white,colbacktitle=keygray,fonttitle=\bfseries,title={$\checkmark$\ Answer Key},arc=3pt,breakable,boxrule=1pt}
\newtcolorbox{recapbox}{colframe=keygray,colback=keybg,coltitle=white,colbacktitle=keygray,fonttitle=\bfseries,title={Question Recap},arc=3pt,breakable,boxrule=1pt}
\newtcolorbox{ideabox}{colframe=titleblue,colback=titleblue!6,coltitle=white,colbacktitle=titleblue,fonttitle=\bfseries,title={The Big Ideas},arc=3pt,breakable,boxrule=1.2pt}
\newtcolorbox{learnbox}{colframe=iagreen,colback=iagreen!5,coltitle=white,colbacktitle=iagreen,fonttitle=\bfseries,title={Concept Essentials},arc=3pt,breakable,boxrule=1.2pt}
\newcommand{\lh}[1]{\par\smallskip{\bfseries\color{iagreen}#1}\par\nobreak\smallskip}

\newcommand{\soln}{\par\smallskip\textit{\textbf{Solution.}}\ }
\newcommand{\workspace}[1]{\par\vspace{3pt}\fcolorbox{black!30}{white}{\begin{minipage}[t][#1][t]{\dimexpr\linewidth-2\fboxsep\relax}\ \end{minipage}}\par}
\newcommand{\sech}[1]{\par\vspace{8pt}{\Large\bfseries\color{iagreen}#1}\par\nobreak\vspace{2pt}{\color{black!20}\hrule height1pt}\vspace{6pt}}

% A compact coordinate plot sized for an example box: \eplot{xmin}{xmax}{ymin}{ymax}{body}
\newcommand{\eplot}[5]{\begin{center}\begin{tikzpicture}
\begin{axis}[width=6.8cm,height=4.4cm,axis lines=middle,xlabel={\tiny$x$},ylabel={\tiny$y$},
grid=both,grid style={gray!16},xmin=#1,xmax=#2,ymin=#3,ymax=#4,
xtick distance=2,ytick distance=2,tick label style={font=\tiny},axis line style={-},samples=120]
#5
\end{axis}\end{tikzpicture}\end{center}}

% A sine-curve plot (degrees on x, 0..360): \splot{ymin}{ymax}{body}
\newcommand{\splot}[3]{\begin{center}\begin{tikzpicture}
\begin{axis}[width=8.6cm,height=4.6cm,axis lines=middle,xlabel={\tiny$x\,(^\circ)$},ylabel={\tiny$y$},
grid=both,grid style={gray!16},xmin=0,xmax=360,ymin=#1,ymax=#2,
xtick distance=90,ytick distance=1,tick label style={font=\tiny},axis line style={-},samples=180]
#3
\end{axis}\end{tikzpicture}\end{center}}

% A small coordinate plot: \plot{xmin}{xmax}{ymin}{ymax}{body}
\newcommand{\plot}[5]{\begin{center}\begin{tikzpicture}
\begin{axis}[width=8.4cm,height=6cm,axis lines=middle,xlabel={\scriptsize$x$},ylabel={\scriptsize$y$},
grid=both,grid style={gray!16},xmin=#1,xmax=#2,ymin=#3,ymax=#4,
xtick distance=2,ytick distance=2,tick label style={font=\tiny},axis line style={-}]
#5
\end{axis}\end{tikzpicture}\end{center}}

% Right triangle (angle theta at bottom-left, right angle at bottom-right):
%   \rtri{drawBase}{drawHeight}{oppLabel}{adjLabel}{hypLabel}
\newcommand{\rtri}[5]{\begin{center}\begin{tikzpicture}[scale=0.85]
\coordinate (A) at (0,0);\coordinate (B) at (#1,0);\coordinate (C) at (#1,#2);
\draw[thick,exblue] (A)--(B)--(C)--cycle;
\draw[black] ($(B)+(-0.32,0)$)--($(B)+(-0.32,0.32)$)--($(B)+(0,0.32)$);
\node[below] at ($(A)!0.5!(B)$){#4};
\node[right] at ($(B)!0.5!(C)$){#3};
\node[above left] at ($(A)!0.5!(C)$){#5};
\node at (0.7,0.22){$\theta$};
\end{tikzpicture}\end{center}}

% General (oblique) triangle with vertex labels and opposite-side labels:
%   \gtri{Alabel}{Blabel}{Clabel}{aSide}{bSide}{cSide}
\newcommand{\gtri}[6]{\begin{center}\begin{tikzpicture}[scale=0.85]
\coordinate (A) at (0,0);\coordinate (B) at (5,0);\coordinate (C) at (1.6,3);
\draw[thick,exblue] (A)--(B)--(C)--cycle;
\node[below left] at (A){#1};\node[below right] at (B){#2};\node[above] at (C){#3};
\node[below] at ($(A)!0.5!(B)$){#6};
\node[right] at ($(B)!0.5!(C)$){#4};
\node[left] at ($(A)!0.5!(C)$){#5};
\end{tikzpicture}\end{center}}

\fancyhead[R]{\footnotesize MCR3U \textperiodcentered{} 3.3}
\begin{document}

\begin{center}
{\LARGE\bfseries\color{titleblue}Worksheet 3.3 --- Transformations of Exponential Functions}\\[2pt]
{\small Grade 11 Functions, University (MCR3U) \textperiodcentered{} Unit 3: Exponential Functions}
\end{center}

\begin{ideabox}
Transform the parent $b^{x}$ with $y=a\,b^{\,k(x-d)}+c$; the shift $c$ moves the asymptote to $y=c$.
\begin{itemize}[leftmargin=*,itemsep=2pt]
  \item $c$ shifts vertically and moves the asymptote to $y=c$.
  \item $d$ shifts horizontally; $a<0$ reflects in the $x$-axis.
  \item Domain stays all reals; range is $y>c$ (or $y<c$ if reflected).
\end{itemize}
\end{ideabox}

\begin{learnbox}
\lh{Shifting the curve}In $y=a\,b^{\,k(x-d)}+c$, the constant $c$ moves the horizontal asymptote to $y=c$ and slides the whole curve up or down.
\lh{Reflections and stretches}$a$ stretches the curve and, if negative, flips it; $d$ shifts it horizontally. The $y$-intercept comes from substituting $x=0$.
\lh{Reading the range}With asymptote $y=c$, the range is $y>c$ when $a>0$ and $y<c$ when $a<0$.
\end{learnbox}

\sech{Worked Examples}

\begin{example}{ex0}{Example 1 --- Vertical shift}
Describe $y=2^{x}+3$.\soln The $+3$ lifts the graph up 3, so the asymptote moves to $y=3$ and the range is $y>3$:\eplot{-4}{3}{-1}{11}{\addplot[red,dashed,domain=-4:3]{3};\addplot[exblue,very thick,domain=-4:3]{2^x+3};}
\end{example}

\begin{example}{ex1}{Example 2 --- Horizontal shift}
Describe $y=2^{x-2}$.\soln $(x-2)$ shifts the graph right 2.
\end{example}

\begin{example}{ex2}{Example 3 --- Reflection}
Describe $y=-2^{x}$.\soln The leading $-$ reflects it in the $x$-axis, so the range becomes $y<0$:\eplot{-3}{3}{-9}{1}{\addplot[exblue,very thick,domain=-3:3]{-(2^x)};}
\end{example}

\begin{example}{ex3}{Example 4 --- Stretch}
Describe $y=3\cdot2^{x}$.\soln A vertical stretch by 3; the $y$-intercept is $3$.
\end{example}

\begin{example}{ex4}{Example 5 --- Range from a shift}
State the range of $y=2^{x}-5$.\soln Asymptote $y=-5$, so the range is $y>-5$.
\end{example}

\begin{example}{ex5}{Example 6 --- Left shift}
Describe $y=2^{x+1}$.\soln $(x+1)$ shifts the graph left 1.
\end{example}

\begin{example}{ex6}{Example 7 --- Reflect in y-axis}
Describe $y=2^{-x}$.\soln The $-x$ reflects it in the $y$-axis, turning growth into decay.
\end{example}

\begin{example}{ex7}{Example 8 --- Asymptote}
State the asymptote of $y=5^{x}+2$.\soln The vertical shift $+2$ gives $y=2$.
\end{example}

\begin{example}{ex8}{Example 9 --- Range with reflection}
State the range of $y=-3^{x}$.\soln Reflected below the axis: $y<0$.
\end{example}

\clearpage
\sech{Your Turn --- Show Your Work}
For each question, show all steps clearly in the space provided.

\begin{qbox}{Question 1}
State the asymptote of $y=2^{x}+6$.\workspace{2.4cm}
\end{qbox}

\begin{qbox}{Question 2}
Describe the shift in $y=3^{x-4}$.\workspace{2.4cm}
\end{qbox}

\begin{qbox}{Question 3}
State the range of $y=2^{x}+2$.\workspace{2.4cm}
\end{qbox}

\begin{qbox}{Question 4}
What does the negative do in $y=-4^{x}$?\workspace{2.4cm}
\end{qbox}

\begin{qbox}{Question 5}
State the asymptote of $y=2^{x}-3$.\workspace{2.4cm}
\end{qbox}

\begin{qbox}{Question 6}
Describe the shift in $y=2^{x+5}$.\workspace{2.4cm}
\end{qbox}

\begin{qbox}{Question 7}
State the range of $y=2^{x}-1$.\workspace{2.4cm}
\end{qbox}

\begin{qbox}{Question 8}
State the domain of $y=3^{x}+7$.\workspace{2.4cm}
\end{qbox}

\begin{qbox}{Question 9}
Describe all transformations of $y=2^{x-1}+4$.\workspace{3cm}
\end{qbox}

\begin{qbox}{Question 10}
State the range of $y=-2^{x}$.\workspace{2.4cm}
\end{qbox}

\begin{qbox}{Question 11}
Write $y=3^{x}$ shifted up $5$.\workspace{2.4cm}
\end{qbox}

\begin{qbox}{Question 12}
Write $y=2^{x}$ shifted left $3$ and down $2$.\workspace{2.4cm}
\end{qbox}

\begin{qbox}{Question 13 --- Challenge}
State the asymptote, domain, and range of $y=-2^{x}+6$.\workspace{3cm}
\end{qbox}

\clearpage
\sech{Question Recap \& Answer Key}

\begin{recapbox}
\begin{multicols}{2}
\small
\begin{enumerate}[leftmargin=*,itemsep=3pt]
  \item State the asymptote of $y=2^{x}+6$.
  \item Describe the shift in $y=3^{x-4}$.
  \item State the range of $y=2^{x}+2$.
  \item What does the negative do in $y=-4^{x}$?
  \item State the asymptote of $y=2^{x}-3$.
  \item Describe the shift in $y=2^{x+5}$.
  \item State the range of $y=2^{x}-1$.
  \item State the domain of $y=3^{x}+7$.
  \item Describe all transformations of $y=2^{x-1}+4$.
  \item State the range of $y=-2^{x}$.
  \item Write $y=3^{x}$ shifted up $5$.
  \item Write $y=2^{x}$ shifted left $3$ and down $2$.
  \item State the asymptote, domain, and range of $y=-2^{x}+6$.
\end{enumerate}
\end{multicols}
\end{recapbox}

\begin{keybox}
\begin{multicols}{2}
\small
\begin{enumerate}[leftmargin=*,itemsep=3pt]
  \item $y=6$
  \item right $4$
  \item $y>2$
  \item reflects over the $x$-axis; range $y<0$
  \item $y=-3$
  \item left $5$
  \item $y>-1$
  \item all real numbers
  \item right $1$, up $4$ (asymptote $y=4$)
  \item $y<0$
  \item $y=3^{x}+5$
  \item $y=2^{x+3}-2$
  \item asymptote $y=6$, domain all reals, range $y<6$
\end{enumerate}
\end{multicols}
\end{keybox}

\end{document}
